\chapter{Media Releases, checksums and \texttt{storeMismatchedFiles}}\label{ch:media-releases}

A Test Schedule cannot run unless a Media Release ships it. As the administrator you do not
install Media Releases through ATE --- they have their own installers --- you put them where the
config says ATE looks, and keep their checksums honest.

\section{Where they go}

ATE searches every directory in \texttt{testDefinitions.directories}, recursively, for
\texttt{.atedef} files. Each Test Definition names a Media Release and lists every Lua file
belonging to it. A directory in that list that does not exist is a warning at each scan, not a
failure:

\begin{lstlisting}[language=ATEShell]
WARNING: Test Definition directory does not exist: /var/lib/ate/definitions
\end{lstlisting}

A Media Release is a directory next to its \texttt{.atedef}:

\begin{lstlisting}[language=ATEShell]
definitions/
  MR-1043/
    release.atedef
    main.tslua
    lib/scpi.lua
\end{lstlisting}

\section{What a Test Definition says}

\texttt{format} must be 1. \texttt{mediaNumber}, \texttt{name} and \texttt{version} are required and
must be non-empty; \texttt{version} must be valid semver 2.0.0. \texttt{files} is a list of
\texttt{path} and \texttt{sha256} pairs, each path ending \texttt{.tslua} (a Test Schedule) or
\texttt{.lua} (a module); \texttt{hidden} is only valid on a schedule, \texttt{module} only on a
module. \texttt{depends} lists other Media Releases by name and a semver range, resolved
transitively, with cycle and ambiguity detection.

A structural defect makes that one definition broken and lists it with its error while the rest of
the scan carries on. The messages are literal, and they are the ones you will see:

\begin{lstlisting}[language=ATEShell]
unknown .atedef format: 2
required 'mediaNumber' is missing or empty
'1.2' is not valid semver 2.0.0
listed file is missing: lib/scpi.lua
duplicate path 'lib/scpi.lua' listed more than once
unknown file extension '.py' for lib/helper.py
depends[0] needs 'name' and 'version'
\end{lstlisting}

A schedule that raises while it is being loaded is reported as broken with its error, and the rest
of its Media Release stays usable.

\section{Checksums, and \texttt{ate definition update}}

Every listed file's SHA-256 is compared with the bytes on disk. A mismatch does not break the
definition; it warns. At scan time the warning goes to stdout; when a Test Run uses the file, one
WARNING line is recorded with the run:

\begin{lstlisting}[language=ATEShell]
checksum mismatch for lib/scpi.lua (expected 3f2a..., got 9b1c...)
\end{lstlisting}

To rewrite the recorded checksums so they match what is actually there --- after a release was
edited, or when a Media Release was shipped without checksums --- run:

\begin{lstlisting}[language=ATEShell]
ate definition update definitions/MR-1043/release.atedef
\end{lstlisting}

It reads the \texttt{.atedef}, computes the SHA-256 of every listed file it can read, and writes
them back in place, filling in any that were absent. The file is only rewritten if at least one
checksum changed; a definition that is already correct is left byte-for-byte as it was. A listed
file that cannot be read keeps its existing \texttt{sha256} and prints a warning rather than
failing the command.

\section{\texttt{database.storeMismatchedFiles}}

This decides what happens to the \emph{bytes} of a mismatched file, not whether the run proceeds:
the run proceeds either way.

\begin{center}
\footnotesize
\begin{tabular}{@{}>{\raggedright\arraybackslash}p{2.6cm}>{\raggedright\arraybackslash}p{12.2cm}@{}}
\toprule
Value & What happens \\
\midrule
\texttt{always} & The raw bytes of every mismatched file are stored with the Test Run. No
  checkbox, no question. \\
\texttt{never} & They are not stored. No checkbox, no question. \\
\texttt{ask} & The Run screen shows a confirmation listing each file as \texttt{expected \dots,
  found \dots} with a "Store mismatched files with this Test Run" checkbox, ticked by default. The
  Operator's answer decides for this run. \\
\bottomrule
\end{tabular}
\end{center}

Headless --- CI, a container, any run with \texttt{QT\_QPA\_PLATFORM=offscreen} --- there is nobody
to click, and the dialog answers itself from the setting: the run starts, and the bytes are stored
only if the setting says \texttt{always}. A run started programmatically with \texttt{ask} and
mismatches pending and no choice given does not guess; it says
\texttt{database.storeMismatchedFiles is 'ask' and this schedule has checksum mismatches: start()
needs store\_mismatched=True or False}.

The stored bytes go into the database, in the \texttt{run\_file} table's \texttt{content} column,
alongside the expected and computed checksums and a mismatch flag. Every Media Release in the
resolved dependency chain of a run is recorded too, with the SHA-256 of its \texttt{.atedef}. That
is why \texttt{always} is a storage decision as much as a debugging one: each mismatched file adds
its whole text to the archive, permanently.

\cando{deploy a Media Release, read a broken Test Definition's error, and decide what
\texttt{storeMismatchedFiles} should say on a Station that runs unattended.}
